\documentclass[12pt]{article}
\usepackage[utf8]{inputenc}
\usepackage[T1]{fontenc}
\usepackage{hyperref}
\title{Tagalog Language Resources for K-12 Curriculum Development}
\author{Compiled by Claude Sonnet 5 (claude-sonnet-5) for Andr\'as Kornai}
\date{}
\begin{document}
\maketitle
\section{Introduction}
Tagalog (ISO 639-3 tgl, Glottocode taga1270) is an Austronesian language of the Philippines with about 45,000,000 speakers. It is used as a language of instruction. The search for this dossier was shallow. It found mostly commercial learner books and one official mother-tongue-based multilingual education (MTB-MLE) item. It found no current DepEd curriculum guides and no Komisyon sa Wikang Filipino materials.
\section{Primary Educational Resources}
\begin{itemize}
\item DepEd (2013), \emph{Mother tongue-based multilingual education: Kagamitan ng mag-aaral Tagalog}, 148 pages, a student book. This is the only official MTB-MLE Tagalog material found. Catalog entry: \url{https://elibrary.lu.edu.ph/students/catalog/2497/}
\item A presentation on MTB-MLE as part of the K to 12 programme. It cites the 2012 implementing guidelines that list Tagalog as a learning area and medium of instruction. \url{https://www.slideshare.net/slideshow/knowing-about-mtbmle/68071062}
\end{itemize}
\section{Linguistic and Reference Materials}
\begin{itemize}
\item Aspillera, \emph{Basic Tagalog} (Tuttle, revised third edition), with a bidirectional dictionary and online audio. It is a learner course for adults and is not school curriculum material. \url{https://www.tuttlepublishing.com/philippines/basic-tagalog}
\item Hippocrene \emph{Tagalog-English/English-Tagalog Standard Dictionary}, about 20,000 entries with a grammar outline. This is a bookseller listing. \url{https://banksquarebooks.com/book/9780781809603}
\end{itemize}
\section{Religious and Community Resources}
No religious or community resources for Tagalog turned up in the searches.
\section{Digital and Audio Resources}
\begin{itemize}
\item University of Iowa Tagalog resources page. It lists free online material, including a Peace Corps manual and grammar handbooks. \url{https://clcl.uiowa.edu/node/266}
\item Online audio accompanies \emph{Basic Tagalog} (see above).
\end{itemize}
\section{Curriculum Development Context}
Tagalog was offered as a learning area and medium of instruction under the 2012 MTB-MLE guidelines. Official DepEd curriculum guides, Filipino-subject textbooks and KWF dictionaries were not found. They should be searched for directly on DepEd and KWF sites.
\end{document}